% ============================================================
%  Chapter 1: Introduction
% ============================================================

\section{Введение}

\subsection{Постановка проблемы}

Понятие \emph{изоморфизма графов} — одно из центральных в дискретной математике.
Два графа изоморфны, если существует биекция между их вершинами, сохраняющая структуру рёбер.
Для классических \emph{невзвешенных} графов изоморфизм хорошо изучен
и в практических задачах решается эффективно~\cite{bowman_graph_iso}.

Однако во множестве современных приложений — биоинформатике, хемоинформатике,
анализе социальных сетей — исходные данные естественно описываются
\emph{полными взвешенными графами} (ПВГ): графами, в которых между каждой парой
вершин имеется ребро с числовым весом.
Теория изоморфизма таких объектов значительно богаче и менее изучена.

\subsection{Актуальность в биоинформатике}

В биоинформатике ПВГ возникают повсеместно:
\begin{itemize}
    \item \textbf{Сети коэкспрессии генов (GCN)} — вершины суть гены,
          вес ребра $(i,j)$ равен коэффициенту корреляции
          между уровнями экспрессии генов $i$ и $j$ по выборке образцов.
    \item \textbf{Матрицы расстояний 3D-структур белков} —
          вершины суть аминокислоты, вес ребра $(i,j)$ ---
          евклидово расстояние между ними в пространстве~\cite{dali1993}.
    \item \textbf{Эволюционные (филогенетические) сети} —
          вершины суть биологические виды или штаммы вирусов,
          вес ребра $(i,j)$ --- число мутаций между их геномами
          (расстояние Хэмминга).
\end{itemize}

Задача, которая возникает в каждом из этих случаев, формально совпадает:
\emph{насколько одинаковы два ПВГ?}
Именно это — проблема изоморфизма (или близости) полных взвешенных графов.

\subsection{Цели и задачи работы}

Целью настоящей работы является разработка математического формализма
для анализа изоморфизма ПВГ и его применение к задаче биомаркерного анализа
в онкологии (сравнение сетей коэкспрессии генов здоровой и опухолевой ткани).

Конкретные задачи:
\begin{enumerate}
    \item Дать строгие определения ПВГ и нескольких понятий изоморфизма.
    \item Сформулировать и доказать основные теоремы (спектральная нижняя оценка,
          NP-трудность, информационно-теоретический предел обнаружения).
    \item Разработать алгоритмы приближённого изоморфизма, пригодные для биологических данных.
    \item Проверить формализм на синтетических данных, моделирующих онкогенез.
\end{enumerate}

\subsection{Структура работы}

Раздел~\ref{sec:theory} вводит математическую теорию ПВГ.
Раздел~\ref{sec:biomodel} переносит её на язык биоинформатики.
Раздел~\ref{sec:experiment} описывает вычислительный эксперимент.
Раздел~\ref{sec:conclusion} подводит итоги.
